% U5-EN — English review, part 1 (adapted from the reviewed Korean U3)
% Purpose: engrXiv/arXiv preprint + IEEE CSM base (free-first strategy)
\documentclass[11pt,a4paper]{article}
\usepackage[T1]{fontenc}
\usepackage{amsmath,amssymb,bm}
\usepackage[margin=25mm]{geometry}
\usepackage{booktabs}
\usepackage{xurl}
\usepackage{hyperref}
\newcommand{\dd}{\mathrm{d}}
\newcommand{\jj}{\mathrm{j}}

\title{Impedance from Telegraph Lines to Robot Contact Control:\\
a Tutorial Review with No Skipped Derivations --- Part 1: Foundations}
\author{Youncheol Jung\\
\small Independent Researcher, Republic of Korea\\
\small ORCID: 0009-0000-5282-881X}
\date{}

\begin{document}
\maketitle

\begin{abstract}
% (abstract inserted from producer's comment block at assembly)
Collaborative robots and humanoids are reaching users who are not control
engineers, yet the language of safe physical interaction --- impedance ---
is usually taught with skipped derivation steps and inconsistent notation.
This tutorial review develops impedance from its historical origin (the
distortion problem of long telegraph lines) through the Laplace transform,
electrical impedance in the RLC circuit, and mechanical impedance in the
mass-spring-damper, with every derivation chain complete and every numeric
example machine-verified and reproducible in an accompanying open-source
MuJoCo teaching laboratory. Part 2 builds robot impedance control and safe
parameter-setting principles on these foundations.
\end{abstract}

% U5-EN — English review manuscript, part 1 (Sections 1–3)
% Derivative of the reviewed Korean manuscript at ../u3_kros/ (adaptation, not literal translation).
% Body only: no preamble. Macros used: \jj, \dd (defined in the main file).
%
% ABSTRACT: With the spread of collaborative robots and humanoids, users who
% are not control specialists increasingly need to configure how a robot
% behaves in contact. The core language of contact behavior is impedance, yet
% most of the literature presents it at an expert level with derivations
% omitted, forcing newcomers to hop between references to reconstruct a
% single concept. Part 1 of this review starts from the historical problem in
% which the language of impedance grew up --- signal distortion on long
% telegraph lines --- and then develops, without skipping derivation steps and
% in a single consistent notation, the definition and differentiation rule of
% the Laplace transform, the electrical impedance of RLC circuits, and the
% mechanical impedance of the mass-spring-damper system. Every numerical
% example can be reproduced in an open educational MuJoCo simulator, and an
% extended manuscript containing all intermediate steps is provided in a
% public repository together with the simulator. Part 2 builds on this
% foundation to treat robot impedance control and the principles of safe
% parameter selection.

\section{Introduction}\label{sec:u5-intro}

Robots have left the factory floor. In homes, hospitals, and small workshops
it is becoming routine for a robot to pick up objects, wipe surfaces, and
insert parts in spaces shared with people \cite{nvidia_physical_ai_2025}.
For such contact tasks, quality and safety are not determined by positioning
accuracy alone; they depend on how the robot has been configured to relate
force to motion --- that is, on its impedance
\cite{hogan1985impedancepart1,abudakka2020variableimpedance}.
The difficulty is that this configuration is not easy even for specialists.
Combine a stiffness with a target offset carelessly and the elastic energy
stored in the virtual spring is released all at once, sending the robot
flying; misunderstand the notion of an equilibrium point and perfectly normal
behavior gets misread as a malfunction.
The author went through exactly this kind of trial and error in front of an
expensive physical robot.

There are two reasons why the existing literature offers little relief.
The first is omitted derivations.
Papers and textbooks skip intermediate steps because of page limits and
assumptions about the reader, and a newcomer ends up wandering across several
references in search of the one step that was left out.
The second is inconsistent notation.
The same physical quantity appears under different symbols in different
sources, so study time is spent matching symbols rather than understanding
concepts.

This review confronts both problems head-on.
It begins at a historical entrance --- showing that impedance is a language
that grew out of a real engineering problem, the signal distortion of
nineteenth-century long-distance telegraph lines --- then derives the Laplace
transform from its definition, and connects electrical (RLC) and mechanical
(mass-spring-damper) impedance in one consistent notation.
For reasons of space, the body of the paper carries the essential
derivations; an extended manuscript (122 pages) containing every intermediate
step, together with an educational MuJoCo simulator that reproduces all
numerical examples, is provided in a public repository.
Part 2 builds on this foundation to cover impedance control, criteria for
choosing between impedance and admittance control, and principles of safe
parameter selection grounded in failure cases.

\section{The Historical Origin of Impedance}\label{sec:u5-history}

Etymologically, the word \emph{impedance} carries the flavor of hindering
progress, but in engineering it means more than simple obstruction: it is a
quantity that expresses both how a system resists an applied flow and how it
stores energy and gives it back.
Why such a concept became necessary is clearest if we return to the problems
of long-distance telegraph and telephone lines in the late nineteenth
century.

\subsection{The long line that broke the resistance-only picture}

For a short DC circuit on a benchtop, one can pack the entire notion of
opposition to flow into a single resistance and explain a great deal with
Ohm's law $V=IR$ alone.
But a telegram is a signal of short pulses switching on and off, and a
telephone call, like a human voice, is a signal that changes continuously;
neither is DC.
On real long cables something strange happened.
Pulses sent over great distances arrived late and rounded off; voices had
their high and low components transmitted differently, so that far away
they arrived smeared and hard to understand.
The sharper the edge of a pulse, the more high-frequency content it needs,
so if the line weakens the high components more, or makes different
frequencies arrive at different times, the distinction between dots and
dashes blurs.
The DC-style intuition --- just lower the wire resistance --- could not
explain this.

The reason is that a long wire is not a simple resistance.
Along a long line there are, distributed with position, a resistance per
unit length $R'~[\Omega/\mathrm{m}]$, an inductance $L'~[\mathrm{H/m}]$, a
capacitance $C'~[\mathrm{F/m}]$, and a conductance $G'~[\mathrm{S/m}]$
representing insulation leakage.
$R'$ and $G'$ are loss elements through which energy disappears as heat and
leakage; $L'$ and $C'$ are storage elements that hold energy briefly in
magnetic and electric fields and return it.
Unlike a lumped-element model that gathers all effects at a single point, a
long line must be viewed as a distributed system in which these small
effects repeat along its length and accumulate.
Writing this idea in terms of the voltage $v(x,t)$ and current $i(x,t)$ at
position $x$ and time $t$ gives the basic form of the telegrapher's
equations,
\begin{align}
  \frac{\partial v}{\partial x} &= -R'i - L'\frac{\partial i}{\partial t}, \\
  \frac{\partial i}{\partial x} &= -G'v - C'\frac{\partial v}{\partial t}.
\end{align}
The first equation says that the voltage changes along the line because of
series resistive loss and the series inductance voltage; the second says
that the current changes because of shunt leakage current and capacitor
charging current.
There is no need to solve these equations here.
It is enough to read them as a sign that a long wire is not one lump of
resistance but an endless chain of small storage and loss elements.

\subsection{Heaviside and the distortionless condition}

Oliver Heaviside encountered this problem first-hand while working as a
telegraph operator, and by applying Maxwell's electromagnetic theory to the
communication-line problem he treated the distortion of long lines
mathematically \cite{heaviside_electrical_papers,donaghyspargo2018heaviside}.
His insight was that the goal is not to send the signal as strongly as
possible, but to make the various frequency components weaken and delay
according to similar rules.
A waveform is a sum of frequency components, so if only some components are
attenuated more or delayed more, the overall shape collapses.
For an idealized line that is uniform and linear, with $R'$, $L'$, $C'$, and
$G'$ nearly independent of frequency, the condition for distortion-free
transmission is
\begin{equation}
  \frac{R'}{L'} = \frac{G'}{C'}.
\end{equation}
$R'/L'$ measures, on the series side, how large the loss is relative to
magnetic-field storage; $G'/C'$ measures, on the shunt side, how large the
leakage is relative to electric-field storage.
Both quantities have units of $\mathrm{s}^{-1}$, so this condition is an
equation matching time scales of the same kind.
Old telephone lines had small leakage $G'$, hence a small $G'/C'$, but
lacked inductance, so $R'/L'$ was comparatively large; adding coils at
regular intervals raised the effective $L'$ and brought the two ratios close
together over the voice band.
Adding inductance --- the very element that opposes changes in current ---
actually improved the signal; this is not an increase in obstruction but a
design that balances loss against storage to reduce waveform distortion.
So the word \emph{impedance} that Heaviside used in the 1880s was not a mere
coinage: it was a working engineer's language for explaining signal
transmission by taking into account inductive and capacitive effects that
resistance alone could not capture
\cite{ethw_heaviside,ptb_impedance_history}.

\subsection{Resistance, reactance, impedance}

For DC, it is often enough to describe the difficulty of flow by a single
resistance.
For AC, however, one must look not only at how much the signal shrinks, but
also at how much the oscillations of voltage and current lead or lag each
other in time.
Picture a practical test: send a pure tone of a single frequency into one
end of a telephone line and measure, at the other end, how much smaller it
has become and how much later it oscillates.
Inductors store energy in magnetic fields and capacitors in electric
fields, and in doing so they shift the phase between voltage and current in
a frequency-dependent way; the frequency-dependent opposition created by
such storage elements is called \emph{reactance}.
Reactance is also measured in ohms ($\Omega$), but unlike resistance it does
not mean that energy is, on average, destroyed as heat.
The AC notion of opposition therefore carries the dissipative component and
the storage component together, written as
\begin{equation}
  Z = R + \jj X.
\end{equation}
The real part $R$ is the resistance that dissipates energy, the imaginary
part $X$ is the reactance that shifts phase, and $\jj$ is the imaginary unit
marking a $90^\circ$ phase difference.
In the sinusoidal steady state the impedances of the three basic elements
are
\begin{equation}
  Z_R = R, \qquad
  Z_L = \jj\omega L, \qquad
  Z_C = \frac{1}{\jj\omega C} = -\frac{\jj}{\omega C},
\end{equation}
where $\omega=2\pi f$ is the angular frequency.
The inductor demands a larger voltage for fast current changes as $\omega$
grows, while the capacitor passes high-frequency currents relatively
easily.
In a series connection the same current flows through all three elements
and the source must supply the sum of the three element voltages, so the
total impedance is the sum of the three impedances,
\begin{equation}
  Z(\jj\omega) = Z_R + Z_L + Z_C
  = R + \jj\left(\omega L - \frac{1}{\omega C}\right).
\end{equation}
The positive reactance of the inductor and the negative reactance of the
capacitor point in opposite phase directions and subtract, and at some
frequency they cancel, leaving the circuit looking as if only the
resistance remained.
This cancellation is the seed of resonance; the detailed development of
resonance itself is treated in the open extended manuscript.
In this way, an AC circuit that would otherwise demand a differential
equation can be handled with the algebraic relation $V=IZ$, an equation
shaped like Ohm's law.
For the detailed derivation of the telegraph-line analysis and the extended
discussion of phasor notation, see the open extended manuscript.

The question the telegraph story leaves behind is how the magnitude and
phase of a time-varying input change as it passes through a system --- and to
write that question down, and to compare electrical circuits and mechanical
systems in the same sentences, we first need the grammar of the Laplace
transform and of linear time-invariant systems.

\section{A Minimal Grammar of the Laplace Transform}\label{sec:u5-laplace}

To turn the question left by the telegraph line --- why and how do a
signal's magnitude and phase change? --- into something computable, two
preparations are needed.
One is an assumption about the system; the other is a transform that turns
differential equations into algebraic ones.
This section lays out that minimal grammar without skipping any derivation
steps.

Start with the assumption.
A system is \emph{linear} when the rule $\mathcal{S}$ that maps an input
$u(t)$ to an output $y(t)$ satisfies
$\mathcal{S}\{a u_1 + b u_2\} = a\,\mathcal{S}\{u_1\} + b\,\mathcal{S}\{u_2\}$
for arbitrary inputs $u_1$, $u_2$ and scalars $a$, $b$.
Double the input and the output doubles; feed in the sum of two inputs and
the result equals the sum of the individual outputs.
A system is \emph{time-invariant} when applying the same input $t_0$ later
merely delays the output by the same $t_0$ without changing its shape,
that is, $\mathcal{S}\{u(t-t_0)\} = y(t-t_0)$.
A real robot is neither perfectly linear nor time-invariant, but for small
motions near a single operating point this linear time-invariant (LTI)
assumption explains a great many phenomena well enough, and it unlocks the
powerful tools of transfer functions and frequency responses.

\subsection{From the definition to the differentiation rule}

The Laplace transform is an integral that summarizes a time function $x(t)$
as a function of a complex variable $s$:
\begin{equation}
  X(s) = \mathcal{L}\{x(t)\}
  = \int_{0}^{\infty} x(t)\,e^{-st}\,\dd t.
  \label{eq:u5-laplace-def}
\end{equation}
In words: multiply $x(t)$ by the steadily shrinking weight $e^{-st}$ and add
everything up from $t=0$ onward.
Each choice of $s$ produces one number, so the entire time waveform is
summarized on the single sheet $X(s)$.

What makes this transform useful is that differentiation turns into
multiplication by $s$ --- and this rule is not a formula to memorize but a
consequence of the definition \eqref{eq:u5-laplace-def} and the product rule
alone.
Since $s$ does not depend on $t$,
\begin{equation}
  \frac{\dd}{\dd t}\left[x(t)\,e^{-st}\right]
  = \dot{x}(t)\,e^{-st} - s\,x(t)\,e^{-st}.
\end{equation}
Integrate both sides from $t=0$ to $\infty$.
The left side integrates a derivative, so it equals the end value minus the
start value: the value of $x(t)e^{-st}$ as $t\to\infty$ minus its value
$x(0)$ at $t=0$.
For values of $s$ at which the response does not blow up, so that
$x(t)e^{-st}\to 0$ as $t\to\infty$, the left side is $0 - x(0) = -x(0)$.
The two integrals on the right side are, by the definition,
$\mathcal{L}\{\dot{x}\}$ and $-sX(s)$ respectively, so
$-x(0) = \mathcal{L}\{\dot{x}\} - sX(s)$, and rearranging gives
\begin{equation}
  \mathcal{L}\{\dot{x}(t)\} = sX(s) - x(0).
  \label{eq:u5-deriv-rule}
\end{equation}
With the initial condition set to zero, $\dot{x}$ can be handled as
$sX(s)$.

Let us run a simple check.
If $x(t)=1$ (a constant), the definition gives
$X(s)=\int_0^\infty e^{-st}\dd t
=\left[-e^{-st}/s\right]_0^\infty
=0-\left(-\tfrac{1}{s}\right)=\tfrac{1}{s}$.
Since $\dot{x}=0$, the left side of rule \eqref{eq:u5-deriv-rule} is
$\mathcal{L}\{0\}=0$, and the right side is
$sX(s)-x(0)=s\cdot(1/s)-1=0$ as well --- they agree.

For a twice-differentiated term, apply the same rule twice.
Treating $\dot{x}$ as a new function and applying
\eqref{eq:u5-deriv-rule} again,
\begin{equation}
  \mathcal{L}\{\ddot{x}\}
  = s\,\mathcal{L}\{\dot{x}\} - \dot{x}(0)
  = s^2 X(s) - s\,x(0) - \dot{x}(0),
\end{equation}
and with zero initial conditions $\ddot{x}$ can be handled as $s^2 X(s)$.

\subsection{From differential equation to transfer function}

Now take the equation of motion of the mass-spring-damper,
$m\ddot{x}+b\dot{x}+kx=f$, and transform it term by term under zero initial
conditions:
$m\ddot{x}\to m s^2 X(s)$, $b\dot{x}\to b s X(s)$, $kx\to kX(s)$, and
$f\to F(s)$, so that
\begin{equation}
  (m s^2 + b s + k)\,X(s) = F(s).
\end{equation}
A differential ledger that had to be updated continuously in time has become
a single page of algebra.
Arranging it as a ratio of output to input gives the \emph{transfer
function},
\begin{equation}
  H(s) = \frac{X(s)}{F(s)} = \frac{1}{m s^2 + b s + k}.
  \label{eq:u5-msd-tf}
\end{equation}
In the denominator, the $ms^2$ term carries the influence of inertia, the
$bs$ term that of damping, and the $k$ term that of the spring's restoring
action.
A transfer function is an input-output relation with the initially stored
energy set to zero (zero-state), and it is worth remembering that the same
physical device presents different transfer functions depending on what is
chosen as input and output.

Here $s$ is not a mere letter of the alphabet but a complex variable that
encodes the character of the time response, usually thought of as
$s=\sigma+\jj\omega$.
$\sigma$ relates to how fast the response grows or decays, and $\omega$ to
how fast it oscillates.
This does not mean that physical inputs are literally complex numbers; it is
a representation that lets decay and oscillation be computed together
inside one symbol.

\subsection{Reading the frequency response}

Feed a stable LTI system the sinusoidal input $u(t)=A\cos(\omega t)$, and
the steady-state output has the same frequency, with only its magnitude and
phase changed:
\begin{equation}
  y_{\mathrm{ss}}(t)
  = A\,|H(\jj\omega)|\cos\!\left(\omega t + \angle H(\jj\omega)\right).
\end{equation}
$|H(\jj\omega)|$ and $\angle H(\jj\omega)$ are the magnitude and angle of
the complex number $H(\jj\omega)$.
Plot a complex number $a+\jj b$ as the point $(a,b)$ in the plane: its
magnitude is the distance from the origin, $\sqrt{a^2+b^2}$, and its angle
is the rotation measured from the positive real axis.
For example, if at some frequency $H(\jj\omega)=1-\jj$, then the magnitude
is $\sqrt{1^2+(-1)^2}=\sqrt{2}$ and the angle is $-45^\circ$, so the output
amplitude is about 1.41 times the input and the phase lags by $45^\circ$.
In other words, the single object $H(\jj\omega)$, obtained by substituting
$s=\jj\omega$, tells us at each frequency how much the system amplifies or
attenuates and how late it responds.
For the relation between poles and response shape, convolution, the
final-value theorem, and other material abbreviated here, see the extended
manuscript.

We are now ready to apply this grammar to its first physical object: in the
next section we transform the equations of the RLC circuit --- resistor,
inductor, capacitor --- by the same procedure and see how electrical
impedance is read from the result.


% U5-EN Sections 4-7 — English adaptation of the reviewed Korean U3
% (sources: u3_kros/u3_s4_electric.tex, u3_s5_mechanical.tex, u3_s6_s7.tex)
\section{Electrical Impedance: the RLC Circuit}
\label{sec:u5-electric}

We now apply the grammar prepared in the previous section to its first
physical embodiment.
A series RLC circuit --- a resistor $R$, an inductor $L$, and a capacitor
$C$ placed on a single current path --- is the smallest electrical system
that can be written as a second-order differential equation.
Starting from the roles of the three elements, we work our way down to the
transfer function and the standard second-order form using the rules of
Section~3, with no derivation steps skipped.

\subsection{The Three Elements: One Dissipation and Two Storages}

The resistor converts electrical energy into heat as current passes through
it.
Under the passive sign convention, in which the terminal where the current
enters is taken as the $+$ of the voltage, we have $v_R = R i$, so the
instantaneous power absorbed by the resistor is $p_R = v_R\,i = R\,i^2$.
Because $i^2$ is always positive regardless of the direction of the
current, the resistor never returns energy: it dissipates continuously.

The inductor stores energy in the magnetic field created by the current,
and satisfies
\begin{equation}
  v_L(t) = L\frac{\dd i(t)}{\dd t}.
  \label{eq:u5-ind-v}
\end{equation}
The faster we try to change the current, the larger the voltage required,
so the inductor resists changes in current.
Substituting \eqref{eq:u5-ind-v} into the instantaneous power and applying
the product rule of differentiation gives
$p_L = v_L\,i = L\,i\frac{\dd i}{\dd t}
= \frac{\dd}{\dd t}\bigl(\tfrac{1}{2}L i^2\bigr)$:
the power flowing into the inductor does not vanish as heat but
accumulates as the rate of change of the magnetic-field energy
$E_L = \frac{1}{2}Li^2$.
During intervals where the current decreases, $p_L<0$ and the stored
energy is returned to the circuit.

The capacitor is the storage element of the opposite kind.
From the separated charge $q = C v_C$, the current is the rate of change
of charge, so
\begin{equation}
  i_C(t) = C\frac{\dd v_C(t)}{\dd t},
  \label{eq:u5-cap-i}
\end{equation}
and since changing the voltage quickly requires moving a large amount of
charge in a short time, the capacitor resists changes in voltage.
Repeating the same product-rule calculation gives
$p_C = v_C\,i_C = C\,v_C\frac{\dd v_C}{\dd t}
= \frac{\dd}{\dd t}\bigl(\tfrac{1}{2}C v_C^2\bigr)$,
so the energy stored in the capacitor is the electric-field energy
$E_C = \frac{1}{2}Cv_C^2$.
Applying the observation from Section~3 that differentiation of a sinusoid
becomes multiplication by $\jj\omega$, the impedances of the three
elements are $Z_R = R$, $Z_L = \jj\omega L$, and $Z_C = 1/(\jj\omega C)$:
the presence or absence of $\jj$ is precisely the difference between
dissipation and storage --- whether the voltage and current are shifted in
phase by $\pm 90^\circ$.

\subsection{The Series RLC Equation and Its Transfer Function}

Let $i(t)$ be the reference current circulating around the loop in one
direction, and let $q(t)$ be the charge on the capacitor plate that this
current enters, so that $i = \dot{q}$.
By Kirchhoff's voltage law, the voltage changes encountered in one full
trip around the loop sum to zero:
$v_{\mathrm{in}} - v_L - v_R - v_C = 0$.
Substituting the element relations $v_L = L\ddot{q}$, $v_R = R\dot{q}$,
and $v_C = q/C$ yields
\begin{equation}
  L\ddot{q}(t) + R\dot{q}(t) + \frac{1}{C}q(t) = v_{\mathrm{in}}(t).
  \label{eq:u5-rlc-ode}
\end{equation}
In this single line, $L\ddot{q}$ is the inductor voltage opposing changes
in current, $R\dot{q}$ is the dissipative voltage, and $q/C$ is a
restoring-type voltage that grows in proportion to the accumulated charge.

We now transform each term, one at a time, using the Laplace
differentiation rule of Section~3 under zero initial conditions:
$L\ddot{q} \to L s^2 Q(s)$, $R\dot{q} \to R s Q(s)$,
$q/C \to Q(s)/C$, and $v_{\mathrm{in}} \to V_{\mathrm{in}}(s)$, so that
\begin{equation}
  \left(L s^2 + R s + \frac{1}{C}\right) Q(s) = V_{\mathrm{in}}(s),
  \qquad
  \frac{Q(s)}{V_{\mathrm{in}}(s)} = \frac{1}{L s^2 + R s + \frac{1}{C}}.
  \label{eq:u5-rlc-tf}
\end{equation}
Choosing the charge as the output is convenient because $i = \dot{q}$ and
$v_C = q/C$ then give the current and the capacitor voltage immediately;
whatever output we pick, the denominator stays the same.

\subsection{The Standard Second-Order Form and Its Two Parameters}

To compare physically different systems on the same scale, second-order
systems are usually written as
\begin{equation}
  \ddot{x} + 2\zeta\omega_n\dot{x} + \omega_n^2 x = \omega_n^2 u,
  \label{eq:u5-std-2nd}
\end{equation}
where $\omega_n$ is the natural frequency, $\zeta$ is the damping ratio,
and $x$ is a temporary name for whichever representative state is being
differentiated twice.
Dividing both sides of \eqref{eq:u5-rlc-ode} by $L$ gives
\begin{equation}
  \ddot{q} + \frac{R}{L}\dot{q} + \frac{1}{LC}q = \frac{1}{L}v_{\mathrm{in}},
\end{equation}
and matching coefficients term by term against \eqref{eq:u5-std-2nd}
yields
\begin{equation}
  2\zeta\omega_n = \frac{R}{L},
  \qquad
  \omega_n^2 = \frac{1}{LC},
  \qquad
  \omega_n^2 u = \frac{1}{L}v_{\mathrm{in}}.
\end{equation}
Substituting $\omega_n^2 = 1/(LC)$ into the last relation gives
$u = C v_{\mathrm{in}}$: the input of the standard form is the voltage
multiplied by the capacitance, a bookkeeping input in units of charge.
This is a good place to see that the standard form changes the scale, not
the physics.
The position term gives $\omega_n = 1/\sqrt{LC}$, and isolating $\zeta$
from the velocity-term comparison,
\begin{equation}
  \zeta
  = \frac{R/L}{2\omega_n}
  = \frac{R}{2L}\sqrt{LC}
  = \frac{R}{2}\sqrt{\frac{LC}{L^2}}
  = \frac{R}{2}\sqrt{\frac{C}{L}}.
\end{equation}
Increasing the resistance increases the damping ratio, while $L$ and $C$
set the rhythm $\omega_n$ at which magnetic-field energy and
electric-field energy trade places.

Let us check this with numbers.
Take $L = 10~\mathrm{mH} = 0.01~\mathrm{H}$ and
$C = 100~\mu\mathrm{F} = 0.0001~\mathrm{F}$, and let $R_c$ denote the
resistance corresponding to critical damping, $\zeta = 1$.
Then
\begin{equation}
  \omega_n = \frac{1}{\sqrt{LC}} = 1000~\mathrm{rad/s},
  \quad
  f_n = \frac{\omega_n}{2\pi} \approx 159~\mathrm{Hz},
  \quad
  R_c = 2\sqrt{\frac{L}{C}} = 20~\Omega.
\end{equation}
If $R$ is smaller than this, the response is underdamped and converges
with oscillation; if larger, it moves toward the overdamped side.
Material trimmed here --- the energy exchange of the free LC oscillation,
series resonance $\omega L = 1/(\omega C)$ and its tuning applications,
and the historical background of the elements --- can be found in the open
expanded manuscript.

In the next section, the same second-order structure reappears as visible
motion: each slot of
$L\ddot{q} + R\dot{q} + q/C = v_{\mathrm{in}}$
is occupied, one for one, by
$m\ddot{x} + b\dot{x} + kx = f$.

\section{Mechanical Impedance: the Mass-Spring-Damper}
\label{sec:u5-mechanical}

Having applied the grammar to its first physical embodiment, we now read
the same second-order structure again as motion we can see.
The mass-spring-damper is the smallest mechanical model for explaining
vibration and impedance.
The mass $m$ is inertia --- a reluctance to change velocity; the stiffness
$k$ is the slope of the restoring force that pulls the displacement back
toward equilibrium; and the damping coefficient $b$ is a resistance that
extracts kinetic energy in proportion to velocity.
Read electrically, the mass corresponds to the inductor, the spring to the
capacitor, and the damper to the resistor.

\subsection{The Equation of Motion: a Ledger of Forces}

Place the equilibrium point at $x=0$ and let $f(t)$ be the external force.
The spring force and the damper force actually acting on the body point in
the directions that reduce the displacement and the velocity, so
$f_s=-kx$ and $f_d=-b\dot{x}$.
Writing Newton's second law \cite{newton1687principia} as
$m\ddot{x}=f+f_s+f_d=f-kx-b\dot{x}$ (the spring force $f_s=-kx$ is
Hooke's law \cite{hooke1678depotentiarestitutiva}) and moving the spring
and damping terms to the left-hand side gives
\begin{equation}
  m\ddot{x}(t)+b\dot{x}(t)+kx(t)=f(t).
  \label{eq:u5m-msd}
\end{equation}
This is less a formula to memorize than a ledger: the burden carried by
the external force $f$ is itemized into the inertia term $m\ddot{x}$, the
damping term $b\dot{x}$, and the spring term $kx$.

\subsection{The Natural Frequency}

We first look at the case with no damping and no external force,
$m\ddot{x}+kx=0$, that is, $\ddot{x}+(k/m)x=0$.
An undamped back-and-forth motion repeats with a fixed shape, so we try
$x(t)=A\cos(\omega t)$.
Differentiating twice gives $\ddot{x}=-\omega^2 x$, so
$-m\omega^2 x+kx=0$; for an oscillation with $x\neq0$ this requires
$\omega^2=k/m$, and therefore $\omega_n=\sqrt{k/m}$.
The units confirm it:
\begin{equation}
  \frac{\mathrm{N/m}}{\mathrm{kg}}
  =
  \frac{\mathrm{kg\,m/s^2}/\mathrm{m}}{\mathrm{kg}}
  =
  \frac{1}{\mathrm{s^2}},
\end{equation}
so taking the square root yields $1/\mathrm{s}$, and treating the radian
as dimensionless we call it rad/s.
The natural frequency is the ratio of the tendency to restore to the
reluctance to move, hence the square-root scaling: quadrupling the
stiffness doubles $\omega_n$, and quadrupling the mass halves it.

\subsection{Standard Form and Coefficient Matching}

Dividing both sides of \eqref{eq:u5m-msd} by $m$ and comparing
coefficients with the standard second-order system
\begin{equation}
  \ddot{x}(t)+2\zeta\omega_n\dot{x}(t)+\omega_n^2 x(t)=\frac{1}{m}f(t)
  \label{eq:u5m-standard}
\end{equation}
gives $\omega_n^2=k/m$ from the $x$ term and $2\zeta\omega_n=b/m$ from the
$\dot{x}$ term, so
\begin{equation}
  \omega_n=\sqrt{\frac{k}{m}}, \qquad \zeta=\frac{b}{2\sqrt{mk}}.
\end{equation}
Here $\omega_n^2$ is the strength of the term that restores position, and
$2\zeta\omega_n$ is the strength of the term that suppresses oscillation
in proportion to velocity.
The damping ratio $\zeta$ is dimensionless because it is $b$ divided by
the critical damping $b_c=2\sqrt{mk}$ introduced below: it packs into a
single number the fact that the same damper acts differently when $m$ and
$k$ differ.

\subsection{Characteristic Roots and the Shape of the Response}

For the free response with $f=0$, substituting $x(t)=e^{st}$ gives the
characteristic equation $ms^2+bs+k=0$, with roots
$s_{1,2}=(-b\pm\sqrt{b^2-4mk})/(2m)$.
The substitution of the standard parameters, which papers usually skip in
one line, we split into two pieces.
Inverting $\zeta=b/(2\sqrt{mk})$ gives $b=2\zeta\sqrt{mk}$, so the front
of the fraction is
\begin{equation}
  -\frac{b}{2m}=-\frac{2\zeta\sqrt{mk}}{2m}=-\zeta\sqrt{\frac{k}{m}}=-\zeta\omega_n,
\end{equation}
and inside the radical, since $b^2=4\zeta^2 mk$,
\begin{equation}
  b^2-4mk=4\zeta^2 mk-4mk=4mk(\zeta^2-1).
\end{equation}
Dividing $\sqrt{b^2-4mk}=2\sqrt{mk}\sqrt{\zeta^2-1}$ by $2m$ gives
$\omega_n\sqrt{\zeta^2-1}$, and assembling the two pieces,
\begin{equation}
  s_{1,2}=-\zeta\omega_n\pm\omega_n\sqrt{\zeta^2-1}.
  \label{eq:u5m-roots}
\end{equation}
Checking with $m=1$, $b=2$, $k=4$: we get $\omega_n=\sqrt{4}=2$ and
$\zeta=2/(2\sqrt{4})=0.5$; the quadratic formula gives
$s_{1,2}=(-2\pm\sqrt{4-16})/2=-1\pm\sqrt{-3}$, and the standard form gives
$-\zeta\omega_n=-1$ and $\omega_n\sqrt{\zeta^2-1}=2\sqrt{-0.75}=\sqrt{-3}$
--- an exact match.

In the underdamped case $0<\zeta<1$, the quantity under the radical is
negative and the roots are complex.
Since $\jj$ is the imaginary unit whose square is $-1$, splitting
$\zeta^2-1=(-1)(1-\zeta^2)$ gives
$\sqrt{\zeta^2-1}=\jj\sqrt{1-\zeta^2}$ (with $1-\zeta^2>0$, the last
radical is an ordinary positive square root), and the characteristic
roots become $s_{1,2}=-\zeta\omega_n\pm\jj\omega_d$.
Here $\omega_d=\omega_n\sqrt{1-\zeta^2}$ is the damped natural frequency,
and the underdamped free response oscillates at $\omega_d$ inside a
shrinking envelope $e^{-\zeta\omega_n t}$:
$x(t)=e^{-\zeta\omega_n t}(A\cos\omega_d t+B\sin\omega_d t)$.

\subsection{Damping Regimes and Critical Damping}

The real part $-\zeta\omega_n$ sets the common rate at which the response
decays, while the square-root term decides whether the roots are real or
complex.
The response oscillates when the discriminant $b^2-4mk$ is negative and
does not when it is positive, so the boundary is $b^2-4mk=0$, that is,
$b=\pm2\sqrt{mk}$.
Since a damping coefficient is physically positive, we take the positive
root: $b_c=2\sqrt{mk}$.
If $\zeta=0$ the oscillation continues indefinitely; if $0<\zeta<1$ the
response oscillates while decaying; if $\zeta=1$ the two roots coincide at
$s=-\omega_n$, giving the boundary response that returns without repeated
oscillation; and if $\zeta>1$ there are two distinct negative real roots,
of which the slower one dominates the overall convergence --- the response
can actually become sluggish.

\subsection{Overshoot and Settling Time}

If a constant force $F_0$ is applied and held, then long afterwards
$\dot{x}=\ddot{x}=0$ and only the spring term remains, so the final value
is $x_{\mathrm{ss}}=F_0/k$.
Overshoot measures how far the response passes this value; for the
standard second-order system with zero initial conditions, a positive
unit-step input, and $0<\zeta<1$, the maximum overshoot ratio is
\begin{equation}
  M_p=\exp\!\left(-\frac{\zeta\pi}{\sqrt{1-\zeta^2}}\right).
  \label{eq:u5m-overshoot}
\end{equation}
When $\zeta$ is small, plenty of energy remains and the response
overshoots strongly: $\zeta=0.2$ predicts about $53\%$.
For $\zeta=0.7$ we get $\sqrt{1-\zeta^2}=\sqrt{0.51}\approx0.714$, the
exponent is $-0.7\pi/0.714\approx-3.08$, and so
$M_p=e^{-3.08}\approx0.046$ --- about $4.6\%$.

The settling time is how long it takes for the response to enter a small
band around the target and stay there.
Treating the amplitude envelope as decaying roughly like
$e^{-\zeta\omega_n t}$ and adopting the $\pm2\%$ criterion --- decay to
$0.02=1/50$ --- we set $e^{-\zeta\omega_n t}\approx0.02$ and take
logarithms of both sides to get
$\zeta\omega_n t\approx\ln 50\approx3.9$.
Since $e^{3.9}\approx50$, this reads as the envelope having shrunk to
about one fiftieth, and from it comes the easy-to-remember approximation
$T_s\approx 4/(\zeta\omega_n)$.
With $\omega_n=10$ and $\zeta=0.7$,
$T_s\approx4/(0.7\cdot10)\approx0.57~\mathrm{s}$.
The approximation is useful roughly for $0.4\lesssim\zeta\lesssim0.8$;
with a 5\% criterion the constant is closer to 3.
The energy ledger that rereads this force-viewpoint vibration as an
exchange between kinetic and elastic potential energy, and the detailed
response plots, are left to the expanded manuscript.

\subsection{Mechanical Impedance}

The same equation of motion goes by different names depending on which
output and which ratio we choose.
Laplace-transforming \eqref{eq:u5m-msd} under zero initial conditions
gives $F(s)=(ms^2+bs+k)X(s)$, and since the velocity is $V_x(s)=sX(s)$,
dividing force by velocity yields the mechanical impedance
\begin{equation}
  Z_m(s)=\frac{F(s)}{sX(s)}=ms+b+\frac{k}{s}.
  \label{eq:u5m-impedance}
\end{equation}
Setting $s=\jj\omega$ gives
$Z_m(\jj\omega)=b+\jj(\omega m-k/\omega)$: the damping term $b$ remains as
a frequency-independent dissipative component, the mass term $\omega m$
grows the faster we shake, and the spring term $k/\omega$ grows the more
slowly we push to change position.
Into the slot where Section~2 previewed the electrical impedance
$Z=R+\jj(\omega L-1/\omega C)$, the damper now steps in for the resistor,
the mass for the inductance, and the spring for the capacitance, in
exactly the same form.
Impedance is therefore not merely a measure of how stiff something is; it
is a force-motion relation that looks spring-like, mass-like, or
damper-like depending on frequency.

This second-order language --- $\omega_n$, $\zeta$, and an impedance that
separates into stiffness, damping, and inertia --- is precisely the
language used in the next section to design the contact behavior of a
robot.

\section{Toward Robot Contact Control}
\label{sec:u5-outlook}

The journey so far compresses into one line.
Impedance is a dynamic relation between force and flow, and that relation
takes the same second-order structure --- two storages and one dissipation
--- in the electrical and the mechanical domains alike.
Robot contact control is the natural next scene for this language.
In impedance control, where a virtual spring and damper are attached to
the robot's fingertip to shape the desired force-motion relation, the
stiffness $k_d$ and the damping ratio $\zeta$ are read on exactly the
scales prepared in this review \cite{hogan1985impedancepart1}.
For example, deciding how much force and how much give to allow at contact
produces a starting value for the stiffness, and the calculation that
switching the stiffness on with a distant target instantly loads
$\frac{1}{2}k_d\delta^2$ of energy is the same spring-energy intuition as
in Section~\ref{sec:u5-mechanical}.

That next scene --- the derivation of impedance control, the criteria for
choosing between it and admittance control, and safe parameter-setting
principles grounded in real accident cases --- is the subject of Part~2 of
this review.
Every numeric example in Parts 1 and 2 can be reproduced in an open
MuJoCo \cite{todorov2012mujoco} teaching simulator.
For instance, the launch accident caused by commanding a large stiffness
without first computing the stored energy can be experienced safely with a
single configuration file in the repository,
\path{configs/lab01_msd/f2_launch_high_energy.yaml}.
The expanded manuscript with all derivations, the simulator code, and the
run guide are provided in the open
repository\footnote{\url{https://github.com/ycpiglet/manipulator-control-tutorial}}.

\section{Conclusion}
\label{sec:u5-conclusion}

Part~1 of this review has organized the language of impedance with no
skipped derivations, in a single notation, and together with its
historical motivation.
Starting from a real problem --- the signal distortion of long telegraph
lines --- we derived the Laplace transform from its definition and
confirmed, by coefficient matching and numeric checks, that the RLC
circuit and the mass-spring-damper share the same second-order structure.
What the reader should take from this journey is not a list of formulas
but the ability to read any second-order system on two scales: the
natural frequency and the damping ratio.
Part~2 carries this ability to the robot's fingertip, addressing the
choice of contact-control scheme and safe parameter setting
\cite{lynchpark2017modernrobotics}.


\bibliographystyle{unsrt}
\bibliography{../../references/refs}
\end{document}
